\documentclass{article}
\usepackage{amsmath}
\usepackage{amssymb}
\usepackage{amsthm}
\usepackage[hidelinks]{hyperref}
\usepackage{url}
\setlength{\emergencystretch}{2em}

\newtheorem{theorem}{Theorem}
\newtheorem{proposition}{Proposition}
\newtheorem{remark}{Remark}
\DeclareMathOperator{\dist}{dist}

\title{Calibrated Selective Constraint Propagation for a Concept-Guided Wall Gate}
\author{Anonymous}
\date{}

\begin{document}
\maketitle

\begin{abstract}
Concept-guided exploration uses an estimated room boundary to restrict later door processing, but a hard geometric threshold can turn a small wall error into a different downstream decision. We isolate the first wall-distance gate of the concept-guided system of Zapata et al. and give it a finite-sample selective certificate. A split-conformal order statistic calibrates a room-level Hausdorff radius from exchangeable, ground-truthed episodes. On the resulting wall-coverage event, every accepted or rejected scan-point decision outside an uncertainty band agrees simultaneously with the decision based on reference walls. This is a marginal guarantee for one interface, not conditional room safety or a certificate for doors, topology, fallback accuracy, or robot motion. The available source papers contain no episode-level wall pairs, so the contribution is a theorem and validation protocol rather than an empirically calibrated system.
% Ledger: 23, 25, 26, 30, 31, 34, 35.
\end{abstract}

\section{Introduction}

Zapata et al. propose a concept-first exploration architecture in which an instantiated room and its walls constrain later door detection \cite{zapata2026concept}. Their implemented detector first retains a scan point $p$ when
\begin{equation}
  \dist(p,W) \leq \delta, \qquad \delta=0.15\,\mathrm{m},
  \label{eq:gate}
\end{equation}
and only then constructs and differentiates a polar scan and pairs peaks. The authors report that high noise can lead to acceptance of an incorrect room model, while revision of accepted models remains future work. They also state that the intended efficiency and robustness gains of hierarchical constraint propagation are not yet quantitatively demonstrated.
% Ledger: 5, 6, 8, 13, 22, 30.

Zhao et al. survey learning-augmented algorithms and emphasise that component guarantees do not compose without explicit relations between interface errors, sensitivity, and downstream objectives \cite{zhao2026learning}. They also require prediction, monitoring, and fallback costs to be accounted for, and discuss split conformal calibration under exchangeability while warning that marginal coverage alone does not control downstream cost.
% Ledger: 6, 8, 12, 21, 26, 30.

This paper connects those observations at one discontinuous interface. The new result is a finite-sample, room-level guarantee: a conformally calibrated Hausdorff radius, combined with a distance-to-set margin, makes all non-abstained decisions of the gate in \eqref{eq:gate} simultaneously agree with a reference-wall gate with marginal probability at least $1-\alpha$. No union bound over scan points is required. The result is deliberately weaker than its possible informal reading: ``simultaneously'' applies only to the first binary gate within a covered new room, and ``probability'' is marginal over an exchangeable new episode.
% Ledger: 23, 25, 26, 30, 31, 34.

Our contribution is therefore the problem-specific coverage-to-interface reduction, not conformal prediction itself and not conformal calibration for robotics. Prior work has calibrated robotic warning systems \cite{luo2021sample}, perception error regions for control \cite{yang2023safe}, and learned perception for navigation \cite{mei2024perceive}. Those papers establish the broader pattern of calibrating an uncertain component and transferring its output to a downstream safety mechanism. We did not locate the specific use of a room-level Hausdorff score to certify every non-abstained decision of a wall-distance gate. This narrower distinction is the novelty claim made here.
% Ledger basis for the narrowed contribution: 25, 26, 30, 34. External-work comparison was established by the searches recorded in search.md.

\section{Setting and source claims}

Let $W^*$ be a nonempty compact reference wall set and $\widehat W$ the nonempty compact wall set returned by a frozen version of the room-fitting pipeline, both in the same finite-dimensional Euclidean space. These assumptions make the distances below finite. Both sets must use the same convention for wall centre-lines or surfaces, observed or complete extent, and reference frame. For a point $p$, write $\dist(p,W)=\inf_{w\in W}\lVert p-w\rVert$. The finite Hausdorff distance is denoted by $d_H$.
% Ledger: 23, 28, 30, 34.

The source system supplies neither a valid Hausdorff-error radius nor episode-level pairs from which one can compute such a radius. In particular, its pose-graph covariance cannot simply be read as a Hausdorff guarantee. This absence motivated replacing a deterministic per-run assumption by a marginal calibration statement.
% Ledger: 20, 22, 24, 35.

\begin{proposition}[First-gate margin]
If $d_H(\widehat W,W^*)\leq\varepsilon$, then, for every $p$,
\begin{equation}
 \left|\dist(p,\widehat W)-\dist(p,W^*)\right|\leq\varepsilon.
 \label{eq:distance}
\end{equation}
Consequently, $\dist(p,\widehat W)+\varepsilon\leq\delta$ certifies acceptance by the reference-wall gate, whereas $\dist(p,\widehat W)-\varepsilon>\delta$ certifies rejection. If neither inequality holds, the point is ambiguous.
% Ledger: 8, 12, 23, 25.
\end{proposition}

\begin{proof}
For any point of one wall set, the Hausdorff bound supplies a point of the other set within $\varepsilon$. Applying the triangle inequality in both directions and taking infima proves \eqref{eq:distance}. If $\dist(p,\widehat W)+\varepsilon\leq\delta$, then \eqref{eq:distance} gives $\dist(p,W^*)\leq\delta$, so the reference gate accepts under its weak inequality. If $\dist(p,\widehat W)-\varepsilon>\delta$, then \eqref{eq:distance} gives $\dist(p,W^*)>\delta$, so the reference gate rejects under the strict complementary condition.
% Ledger: 8, 12, 34.
\end{proof}

Without a margin, the binary gate is not globally Lipschitz with respect to Hausdorff perturbations: an arbitrarily small wall translation can move a point across the threshold. This statement concerns membership at the first gate. It does not by itself show that a particular door or topology edge changes.
% Ledger: 8, 12. The final sentence makes explicit the weaker supported reading in entries 9 and 13.

\section{Calibrated selective certificate}

Fix a positive integer $n$ and a target miscoverage level $0<\alpha<1$. Freeze the complete room fitter independently of calibration. Let
$(\widehat W_i,W_i^*)$, $i=1,\ldots,n$, be calibration episodes and define
\begin{equation}
 R_i=d_H(\widehat W_i,W_i^*), \qquad
 k=\left\lceil(n+1)(1-\alpha)\right\rceil.
 \label{eq:scores}
\end{equation}
Let $\varepsilon_\alpha$ be the $k$th order statistic of $R_1,\ldots,R_n$, with $\varepsilon_\alpha=+\infty$ when $k>n$. Calibration and the new episode are assumed exchangeable at the episode level.
% Ledger: 23, 25, 26, 30, 34.

\begin{theorem}[Marginal simultaneous correctness of the first gate]
Under the assumptions above,
\begin{equation}
 \Pr\!\left\{d_H(\widehat W_{\mathrm{new}},W^*_{\mathrm{new}})
 \leq\varepsilon_\alpha\right\}\geq 1-\alpha.
 \label{eq:coverage}
\end{equation}
For the new room, certify acceptance when
$\dist(p,\widehat W_{\mathrm{new}})+\varepsilon_\alpha\leq\delta$, certify rejection when
$\dist(p,\widehat W_{\mathrm{new}})-\varepsilon_\alpha>\delta$, and abstain otherwise. Then
\begin{equation}
 \begin{aligned}
 \Pr\!\{&\text{all non-abstained first-gate decisions agree}\\
         &\text{with the reference-wall gate}\}\geq 1-\alpha.
 \end{aligned}
 \label{eq:transfer}
\end{equation}
% Ledger: 23, 25, 26, 30, 34.
\end{theorem}

\begin{proof}
Break ties among the $n+1$ exchangeable scores using independent continuous random variables. The resulting randomized rank of the new score is uniform on $\{1,\ldots,n+1\}$, so the choice of $k$ gives the conformal coverage bound. Replacing randomized tie-breaking by the non-randomized weak event in \eqref{eq:coverage} can only make coverage conservative. The infinity convention covers unattainable finite quantiles. On the single event in \eqref{eq:coverage}, Proposition 1 applies to every scan point. Each non-abstained decision consequently matches its reference-wall counterpart. Because one room-level event covers all points, there is no pointwise union bound.
% Ledger: 23, 25, 26, 34.
\end{proof}

The theorem is selective: it permits abstention, and it does not bound how often abstention occurs. Its probability is marginal across exchangeable rooms or trajectory clusters, rather than conditional on a particular room, geometry, or noise regime. Correlated scan points or frames must not be counted as independent calibration episodes.
% Ledger: 24, 26, 28, 31, 35.

\section{Fallback and cost interpretation}

An ambiguous point cannot safely be routed in isolation. Filtering changes the samples used to construct the polar scan, and hence can alter angular adjacency, derivatives, peaks, and peak pairs. A declared fallback must therefore cover the whole scan or a dependency region equipped with a fixed-grid, interpolation, and boundary contract. This is an interface requirement, not a proved fallback algorithm.
% Ledger: 13, 14, 15, 18, 30, 31.

For completeness, we state a separate compute-only observation in a self-contained model. Let $I$ be a fixed perception instance, namely the current observable scene, and let $h$ be the room hypothesis used as structural advice. Algorithm $A$ searches the wall-constrained regions specified by $h$; algorithm $B$ searches the declared instance class independently of $h$ and is complete on that class. Let $C_A(I,h)$ and $C_B(I)$ be the amounts of one common divisible compute resource that the respective algorithms require to return a result passing the same independent validity check; $C_A$ may be infinite, while completeness makes $C_B$ finite. Both algorithms are assumed preemptible and to preserve their state between resource slices. The mixed execution allocates resource shares $1-\lambda$ and $\lambda$, where $0<\lambda<1$, to $A$ and $B$, respectively, and stops when either result passes the validity check. If $C_{\mathrm{mix}}(I,h)$ denotes the total common-resource work consumed by that execution before it stops, then
% Ledger: 3, 4, 10, 18.
\begin{equation}
 C_{\mathrm{mix}}(I,h)\leq
 \min\left\{\frac{C_A(I,h)}{1-\lambda},
             \frac{C_B(I)}{\lambda}\right\}.
 \label{eq:dovetail}
\end{equation}
This is the standard time-sharing argument. Neither source paper supplies the required complete baseline or independent validator, so \eqref{eq:dovetail} is not an implemented-system guarantee and says nothing about policy-dependent motion.
% Ledger: 3, 4, 10, 15, 18, 31.

\section{Validation protocol}

The smallest empirical test fixes a deployment population and wall-set convention, freezes the fitter, and collects disjoint ground-truthed calibration and test rooms in the source system's simulator. It should report room-level coverage, $\varepsilon_\alpha/\delta$, the fraction of scan dependency regions sent to fallback, and first-gate disagreement against reference walls. Door recall, topology error, compute, and motion should be reported separately, without treating them as consequences of the theorem.
% Ledger: 24, 26, 28, 30, 31, 35.

The study should include the source paper's planned noise variation and accepted-wrong-room failure modes, as well as near-threshold layouts. A held-out change of building, sensor, or noise distribution should be labelled as distribution shift rather than covered by \eqref{eq:coverage}. If $\varepsilon_\alpha\geq0.15\,\mathrm{m}$ or dependency expansion yields nearly total fallback, the certificate can remain valid while being operationally uninformative; that outcome should be reported as negative.
% Ledger: 13, 24, 28, 30, 31, 35.

\section{Limitations and open questions}

First, exchangeability is a substantive sampling assumption. Distribution shift and dependent adaptive sampling invalidate the stated rank guarantee, and the theorem supplies no conditional coverage for rare room types.
% Ledger: 24, 26, 31, 34.

Second, wall-set conventions and frames must match. A score between a centre-line estimate and a surface reference, or between observed and complete extents, does not instantiate the theorem's intended quantity.
% Ledger: 28, 31, 34.

The source system is restricted to structured indoor environments: it models rectangular rooms with orthogonal walls under a Manhattan-world assumption, and its detector thresholds were tuned empirically for its particular setup. The theorem's distance-to-set statement is geometric, but the proposed experiment, calibration population, and any operational conclusion necessarily inherit the chosen detector and deployment population. They should not be read as evidence for non-rectangular or otherwise unrepresented environments.
% Ledger: 15, 24, 26, 30, 31, 35.

Third, the certificate stops at the first wall-distance gate. It does not certify polar resampling, derivative thresholds, peak pairing, width tests, segment tests, final doors, graph topology, fallback accuracy, or active robot motion. The downstream dependency can also turn a narrow pointwise ambiguity band into broad regional fallback.
% Ledger: 9, 13, 14, 18, 24, 28, 31, 35.

Fourth, the source material contains no calibration data, so neither the radius, achieved coverage, nor abstention cost is known. Practical non-vacuity remains an empirical question.
% Ledger: 28, 35.

Open questions are whether a defensible deployment population can be specified, whether useful coverage and regional abstention can coexist at $\delta=0.15\,\mathrm{m}$, and whether later detector stages admit compatible interface margins. A separate open problem is to specify and validate a complete fallback before extending the compute observation to active sensing.
% Ledger: 15, 26, 28, 31, 35.

\section{Conclusion}

The paired reading yields a precise but limited result. Split-conformal room-level coverage of a Hausdorff wall error transfers, through the distance-to-set inequality, to simultaneous correctness of every non-abstained decision at a concept-guided detector's first wall gate. Existing robotics work already establishes conformal calibration as a safety tool; the contribution here is this interface-specific transfer and its explicit abstention contract. Whether it is useful cannot be decided from the two source papers and requires the stated calibration experiment.
% Ledger: 23, 25, 26, 30, 34, 35.

\bibliographystyle{plain}
\bibliography{references}

\end{document}
